\chapter{The Albanese universal property of \(\Pic^0_{C/k}\)}
\label{chap:Albanese_AlbaneseUP}
% archon:covers AlgebraicJacobian/Albanese/AlbaneseUP.lean

% STRATEGY NOTE (iter-174 / refined iter-175). This chapter is the A.4.d
% sub-row of Route~A (the positive-genus arm of
% \texttt{nonempty\_jacobianWitness}). It is the join of two preceding rows:
% A.3 (the identity component \(\Pic^0_{C/k}\) of the Picard scheme together
% with its degree map; the relevant Picard-scheme representability content
% is recorded in \cref{chap:Picard_FGAPicRepresentability}, A.2.c), and
% A.4.c (\cref{chap:Albanese_Thm32RationalMapExtension}, Milne's
% Theorem~3.2 on rational maps from a nonsingular variety to an abelian
% variety). The output is Milne's Proposition~6.1 of Chapter~III~\S 6 of
% \emph{Abelian Varieties}: the Abel--Jacobi morphism
% \(\iota_{P_0} \colon C \to \Pic^0_{C/k}\) at a marked \(\bar k\)-point
% \(P_0 \in C(\bar k)\) is the \emph{Albanese morphism} of the pointed curve
% \((C, P_0)\). Any pointed morphism \(\varphi \colon C \to A\) into an abelian
% variety with \(\varphi(P_0) = \eta_A\) factors uniquely as
% \(\varphi = \psi \circ \iota_{P_0}\) for a homomorphism of group schemes
% \(\psi \colon \Pic^0_{C/k} \to A\).
%
% ROUTE DECISION (iter-174, refined iter-175). Two routes were on the table:
%   (i) MODULI / AUTODUALITY route. Pull the Poincar\'e bundle of \(A\) back
%       along \(\varphi \times \mathrm{id}_{A^\vee}\), classify the resulting
%       trivialised relative degree-zero line bundle by the Yoneda map
%       \(A^\vee \to J\), dualise via the autoduality \(J \cong J^\vee\)
%       (Milne~III~6.6) and the principal polarisation. This route depends
%       on the autoduality of \(J\), which gates on the \emph{theorem of the
%       cube} -- EXCISED from the project in iter-163.
%  (ii) SYMMETRIC-POWER route. Define $\widetilde\psi \colon \mathrm{Sym}^g
%       C \to A\( from the symmetric assignment \)(P_1,\ldots,P_g) \mapsto
%       \varphi(P_1) + \cdots + \varphi(P_g)$, descend through the
%       birational morphism \(\mathrm{Sym}^g C \dashrightarrow J\), and apply
%       Milne's Theorem~3.2 (\cref{thm:rational_map_to_av_extends}) to
%       promote the rational map \(J \dashrightarrow A\) to a regular
%       morphism. This is Milne's own proof of Proposition~6.1; it uses
%       \emph{only} proven \cref{thm:rigidity_lemma} (Milne~I~1.1) and
%       \cref{thm:rational_map_to_av_extends} (Milne~I~3.2), and avoids the
%       cube entirely.
% The project commits to Route~(ii). The moduli sub-lemmas from the iter-174
% draft (Poincar\'e-bundle pullback, moduli classifier of a trivialised line
% bundle) are retired; an ``Alternative route history'' \% NOTE block at the
% end of this chapter records the moduli wording for reference.
%
% This is the final wiring of the positive-genus Albanese-witness
% construction: the conclusion supplies, for each \(\bar k\)-rational point
% \(P_0 \in C(\bar k)\), the universal-factorisation field
% \(\texttt{isAlbaneseFor}\,P_0\) of an Albanese witness for \((C, P_0)\) in the
% sense of \cref{def:IsAlbanese} of \cref{chap:Jacobian}; together with the
% group-object / smooth / proper / geometrically-irreducible structure on
% \(\Pic^0_{C/k}\) obtained at A.3 and the Abel--Jacobi morphism of
% \cref{lem:abel_jacobi_morphism} below, this yields a positive-genus
% \texttt{JacobianWitness} for \(C/\bar k\) which then descends back to \(k\) by
% Galois descent (\cref{chap:Jacobian}, sub-step C.2.f).
%
% The Lean target file is \texttt{AlgebraicJacobian/Albanese/AlbaneseUP.lean}
% (not yet created); this chapter is the prover-ready specification.

\section{Setup and statement}
\label{sec:albaneseup_setup}

Fix an algebraically closed field \(\bar k\) and a smooth proper geometrically
irreducible curve \(C\) over \(\bar k\) of genus \(g = g(C) > 0\) (the positive-genus
hypothesis is essential: it is precisely the standing hypothesis of Milne's
\S III.6). Fix moreover a \(\bar k\)-rational point \(P_0 \in C(\bar k)\) to serve
as the basepoint of the Albanese construction.

By the A.3 row of Route~A (the identity-component construction applied to the
representable Picard scheme of
\cref{def:pic_scheme} in \cref{chap:Picard_FGAPicRepresentability}) the
identity component \(\Pic^0_{C/k} \subseteq \Pic_{C/k}\) is a smooth proper
geometrically irreducible \(\bar k\)-group scheme of dimension \(g\); we write
\(J := \Pic^0_{C/k}\) and let \(\eta_J \colon \Spec \bar k \to J\) denote its
identity section, in the project's convention the chosen \(\mathbf 1\)-pointed
morphism that represents the unit element of the underlying group object on
\(J\).

The universal \(T\)-points of \(J\) are degree-\(0\) line bundles on \(C \times_{\bar k}
T\), taken modulo line bundles pulled back from \(T\) (\cref{chap:Picard_LineBundlePullback};
the identity component then refines this to the degree-zero subfunctor on the
\'etale sheafification). This moduli interpretation is what underlies the
Abel--Jacobi morphism \(\iota_{P_0} \colon C \to J\) of
\cref{lem:abel_jacobi_morphism}. The universal homomorphism \(\psi \colon J
\to A\) corresponding to a pointed morphism \(\varphi \colon C \to A\) is
constructed by Milne's symmetric-power route: the symmetric assignment
\((P_1, \ldots, P_g) \mapsto \varphi(P_1) + \cdots + \varphi(P_g)\) defines a
regular morphism \(\widetilde \psi \colon \mathrm{Sym}^g C \to A\), the
birational morphism \(\mathrm{Sym}^g C \dashrightarrow J\) inverts to a
rational map \(J \dashrightarrow A\), and the rational-map extension theorem
\cref{thm:rational_map_to_av_extends} promotes that rational map to the
desired regular morphism \(\psi\).

The goal of the chapter is to establish:

\begin{theorem}
[Albanese universal property of \(\Pic^0_{C/k}\) -- Milne III~\S 6, Proposition~6.1]
  \label{thm:albanese_universal_property}
  \lean{AlgebraicGeometry.Pic0.albanese_universal_property}
  \uses{lem:abel_jacobi_morphism, lem:symmetric_product_av_map, lem:symmetric_product_to_jacobian, lem:descent_through_birational_sigma, lem:albanese_eq_iff_symmetricPower_eq, thm:rational_map_to_av_extends, thm:rigidity_lemma}

  % SOURCE: [Milne, Abelian Varieties], Proposition~6.1, \S III.6, doc p.~104
  % (read from references/abelian-varieties.pdf, PDF page 110).
  % SOURCE QUOTE:
  % "PROPOSITION 6.1. Let \(P\) be a \(k\)-rational point on \(C\).  The map
  %  \(f^P\colon C\,\to\,J\) has the following universal property: for any
  %  map \(\varphi\colon C\to A\) from \(C\) into an abelian variety sending
  %  \(P\) to \(0\), there is a unique homomorphism \(\psi\colon J\to A\) such
  %  that \(\varphi=\psi\circ f^P\)."
  \textit{Source: Milne, \emph{Abelian Varieties}, Proposition~6.1, \S III.6, p.~104.}

  Let \(C\) be a complete nonsingular curve of genus \(g > 0\) over an algebraically
  closed field \(\bar k\), let \(P_0 \in C(\bar k)\) be a \(\bar k\)-rational point,
  and write \(J := \Pic^0_{C/k}\), \(\iota_{P_0} \colon C \to J\) for the
  Abel--Jacobi morphism of \cref{lem:abel_jacobi_morphism}. Then for every
  abelian variety \(A\) over \(\bar k\) and every pointed morphism
  \(\varphi \colon C \to A\) with \(\varphi(P_0) = \eta_A\), there exists a unique
  homomorphism of group schemes \(\psi \colon J \to A\) such that
  \(\varphi = \psi \circ \iota_{P_0}\).
\end{theorem}

The rest of this chapter develops the sub-lemmas
(\cref{lem:abel_jacobi_morphism}, \cref{def:symmetric_power_curve},
\cref{lem:symmetric_product_av_map}, \cref{lem:symmetric_product_to_jacobian},
and \cref{lem:descent_through_birational_sigma}) that the proof rests on,
then gives the proof.

\section{File-internal placeholder for \(\Pic^0_{C/k}\)}
\label{sec:albaneseup_pic0_placeholder}

The A.3 row of Route~A (the identity-component refinement of the representable
Picard scheme together with its degree map) supplies the Jacobian scheme
\(J := \Pic^0_{C/k}\) as a smooth proper geometrically irreducible
\(\bar k\)-group scheme. Pending a dedicated A.3 chapter, this section records
the four structural attributes of \(J\) as a single bundled placeholder, from
which the underlying scheme and its abelian-variety instances are projected.
Once the A.3 chapter materialises, this placeholder collapses to a re-export.

\begin{definition}\leanok
[Pic\(^0\) structural bundle]
  \label{def:pic0_bundle}
  \lean{AlgebraicGeometry.Pic0.Bundle}
  \uses{def:pic_scheme}

  A \emph{Pic\(^0\) bundle} for a smooth proper geometrically irreducible
  curve \(C\) over \(\bar k\) is the datum of the underlying \(\bar k\)-scheme
  \(\Pic^0_{C/k}\) together with its four abelian-variety structural
  attributes obtained at the A.3 row of Route~A: a group-object structure,
  properness over \(\bar k\), smoothness over \(\bar k\), and geometric
  irreducibility. It packages, as a single object, the output of the
  identity-component refinement of the representable Picard scheme of
  \cref{def:pic_scheme}.
\end{definition}

\begin{definition}\leanok
[Pic\(^0\) bundle carrier]
  \label{def:pic0_bundle_carrier}
  \lean{AlgebraicGeometry.Pic0.bundle}
  \uses{def:pic0_bundle}

  For each smooth proper geometrically irreducible curve \(C\) over \(\bar k\)
  there is a Pic\(^0\) bundle (\cref{def:pic0_bundle}), supplied by the A.3
  row: the identity component \(\Pic^0_{C/k}\) of the representable Picard
  scheme, equipped with its group-object, proper, smooth, and geometrically
  irreducible structure.
\end{definition}

\begin{definition}\leanok
[The Jacobian scheme \(\Pic^0_{C/k}\)]
  \label{def:jacobian_scheme}
  \lean{AlgebraicGeometry.Pic0.jacobianScheme}
  \uses{def:pic0_bundle_carrier}

  The \emph{Jacobian scheme} \(J := \Pic^0_{C/k}\) of \(C\) is the underlying
  \(\bar k\)-scheme of the Pic\(^0\) bundle of
  \cref{def:pic0_bundle_carrier}.
\end{definition}

\begin{lemma}\leanok
[Group-object structure on \(\Pic^0_{C/k}\)]
  \label{lem:jacobian_scheme_grpObj}
  \lean{AlgebraicGeometry.Pic0.jacobianScheme_grpObj}
  \uses{def:jacobian_scheme}

  The Jacobian scheme \(J = \Pic^0_{C/k}\) (\cref{def:jacobian_scheme})
  carries a group-object structure, read off from the corresponding field of
  its Pic\(^0\) bundle.
\end{lemma}

\begin{proof} Proved directly in Lean. \end{proof}

\begin{lemma}\leanok
[Properness of \(\Pic^0_{C/k}\)]
  \label{lem:jacobian_scheme_isProper}
  \lean{AlgebraicGeometry.Pic0.jacobianScheme_isProper}
  \uses{def:jacobian_scheme}

  The Jacobian scheme \(J = \Pic^0_{C/k}\) (\cref{def:jacobian_scheme}) is
  proper over \(\bar k\), read off from the corresponding field of its
  Pic\(^0\) bundle.
\end{lemma}

\begin{proof} Proved directly in Lean. \end{proof}

\begin{lemma}\leanok
[Smoothness of \(\Pic^0_{C/k}\)]
  \label{lem:jacobian_scheme_smooth}
  \lean{AlgebraicGeometry.Pic0.jacobianScheme_smooth}
  \uses{def:jacobian_scheme}

  The Jacobian scheme \(J = \Pic^0_{C/k}\) (\cref{def:jacobian_scheme}) is
  smooth over \(\bar k\), read off from the corresponding field of its
  Pic\(^0\) bundle.
\end{lemma}

\begin{proof} Proved directly in Lean. \end{proof}

\begin{lemma}\leanok
[Geometric irreducibility of \(\Pic^0_{C/k}\)]
  \label{lem:jacobian_scheme_geomIrred}
  \lean{AlgebraicGeometry.Pic0.jacobianScheme_geomIrred}
  \uses{def:jacobian_scheme}

  The Jacobian scheme \(J = \Pic^0_{C/k}\) (\cref{def:jacobian_scheme}) is
  geometrically irreducible over \(\bar k\), read off from the corresponding
  field of its Pic\(^0\) bundle.
\end{lemma}

\begin{proof} Proved directly in Lean. \end{proof}

\section{The Abel--Jacobi morphism at a marked point}
\label{sec:albaneseup_abel_jacobi}

The first sub-lemma packages the construction and the elementary regularity of
the Abel--Jacobi morphism \(\iota_{P_0}\) at the basepoint \(P_0\). It is the
content gated on A.3 (existence of \(\Pic^0_{C/k}\) as a representable group
scheme).

\begin{lemma}\leanok
[Abel--Jacobi morphism at a basepoint]
  \label{lem:abel_jacobi_morphism}
  \lean{AlgebraicGeometry.Pic0.abelJacobi}
  \uses{def:pic_scheme, def:jacobian_scheme}

  % SOURCE: [Milne, Abelian Varieties], \S III.6, the canonical map
  % \(f^P\colon C\to J\) (introduction to Proposition~6.1), doc p.~104
  % (read from references/abelian-varieties.pdf, PDF page 110).
  % SOURCE QUOTE:
  % "PROPOSITION 6.1. Let \(P\) be a \(k\)-rational point on \(C\).  The map
  %  \(f^P\colon C\to J\) has the following universal property: for any map
  %  \(\varphi\colon C\to A\) from \(C\) into an abelian variety sending \(P\)
  %  to \(0\), there is a unique homomorphism \(\psi\colon J\to A\) such that
  %  \(\varphi=\psi\circ f^P\)."
  % SOURCE: [Milne, Abelian Varieties], \S III.6, Summary 6.11, doc p.~107
  % (read from references/abelian-varieties.pdf, PDF page 113).
  % SOURCE QUOTE:
  % "SUMMARY 6.11.  Between \((C,P)\) and itself, there is a divisorial
  %  correspondence \(\mathcal L^P = \mathcal L(\Delta-\{P\}\times C-C\times\{P\})\).
  %    Between \((C,P)\) and \(J\) there is the divisorial correspondence
  %  \(\mathcal M^P\); for any divisorial correspondence \(\mathcal L\) between
  %  \((C,P)\) and a pointed \(k\)-scheme \((T,t)\), there is a unique morphism
  %  of pointed \(k\)-schemes \(\varphi\colon T\to J\) such that
  %  \((1\times \varphi)^*\mathcal M^P\approx \mathcal L\).  In particular,
  %  there is a unique map \(f^P\colon C\to J\) such that
  %  \((1\times f^P)^*\mathcal M^P\approx \mathcal L^P\) and \(f(P)=0\)."
  \textit{Source: Milne, \emph{Abelian Varieties}, Proposition~6.1 and
  Summary~6.11, \S III.6.}

  Let \(C\) be a complete nonsingular curve of genus \(g > 0\) over \(\bar k\), let
  \(P_0 \in C(\bar k)\) be a basepoint, and write \(J := \Pic^0_{C/k}\). There is a
  unique regular morphism of \(\bar k\)-schemes
  \[
    \iota_{P_0} \;\colon\; C \;\longrightarrow\; J, \qquad
    Q \;\longmapsto\; \bigl[\mathcal O_C(Q - P_0)\bigr],
  \]
  sending the basepoint \(P_0\) to the identity element \(\eta_J \in J(\bar k)\)
  and characterised on \(T\)-points by the formula
  \((1 \times \iota_{P_0})^*\mathcal M^{P_0} \approx \mathcal L^{P_0}\), where
  \(\mathcal L^{P_0} := \mathcal O_{C \times C}\bigl(\Delta - \{P_0\} \times C -
  C \times \{P_0\}\bigr)\) is the rigidified diagonal correspondence on
  \(C \times C\) and \(\mathcal M^{P_0}\) is the Poincar\'e--Picard correspondence
  on \(C \times J\).
\end{lemma}

% SOURCE QUOTE PROOF: (no separate verbatim source; the existence of
% \(f^P\) is part of Milne's standing setup at the head of \S III.6 -- ``\(J\)
% will be its Jacobian variety'' -- and is the moduli-of-degree-0-line-bundles
% interpretation of \(\Pic^0_{C/k}\) packaged into Summary~6.11.)
\begin{proof}
  
  By A.3, \(J := \Pic^0_{C/k}\) is the moduli scheme of degree-zero line bundles
  on \(C\); concretely, for every \(\bar k\)-scheme \(T\) its \(T\)-points are
  isomorphism classes of line bundles \(\mathcal L\) on \(C \times_{\bar k} T\)
  whose fibre degree at every geometric point of \(T\) is zero, taken modulo line
  bundles pulled back from \(T\) (the universal property of the Picard scheme
  packaged in \cref{def:pic_scheme}; the identity component imposes the
  degree-zero condition).

  Apply this to \(T := C\) and the line bundle
  \[
    \mathcal L^{P_0} \;:=\; \mathcal O_{C \times_{\bar k} C}\bigl(\Delta -
      \{P_0\} \times C - C \times \{P_0\}\bigr),
  \]
  the rigidified diagonal correspondence on \(C \times_{\bar k} C\). Its fibre
  over a \(\bar k\)-point \(Q \in C(\bar k)\) along the second projection is the
  degree-zero line bundle \(\mathcal O_C(Q - P_0)\) on \(C\), so \(\mathcal L^{P_0}\)
  is a relative degree-zero line bundle, hence a \(C\)-point of \(J\). The
  representability of \(J\) at \(T = C\) produces the unique morphism
  \[
    \iota_{P_0} \;\colon\; C \;\longrightarrow\; J
  \]
  such that \((1 \times \iota_{P_0})^*\mathcal M^{P_0} \approx \mathcal L^{P_0}\)
  for the Poincar\'e--Picard correspondence \(\mathcal M^{P_0}\) on \(C \times J\).

  At the basepoint, restriction of \(\mathcal L^{P_0}\) to \(C \times \{P_0\}\)
  gives \(\mathcal O_C(P_0 - P_0) \cong \mathcal O_C\), the trivial line bundle,
  whose moduli class in \(J(\bar k)\) is the identity element \(\eta_J\). Hence
  \(\iota_{P_0}(P_0) = \eta_J\). Uniqueness of \(\iota_{P_0}\) is the moduli
  uniqueness clause of representability: any other morphism \(\iota'_{P_0}\)
  satisfying \((1 \times \iota'_{P_0})^*\mathcal M^{P_0} \approx \mathcal L^{P_0}\)
  and the pointed condition agrees with \(\iota_{P_0}\) on \(T\)-points for every
  \(T\), hence equals \(\iota_{P_0}\) by the Yoneda embedding.
\end{proof}

\section{The \(g\)-th symmetric power \(\mathrm{Sym}^g C\)}
\label{sec:albaneseup_sym_g}

The Albanese universal property is built from \(C\) by taking the
\(g\)-th \emph{symmetric power} \(\mathrm{Sym}^g C\) -- the quotient of the
\(g\)-fold Cartesian power \(C^g\) by the natural action of the symmetric group
\(S_g\). We collect the definition here as the first sub-lemma block.

\begin{definition}\leanok
[Symmetric power of a curve]
  \label{def:symmetric_power_curve}
  \lean{AlgebraicGeometry.Pic0.SymmetricPower}
  % SOURCE: [Milne, Abelian Varieties], Proposition~3.1, \S III.3,
  % doc p.~94 (read from references/abelian-varieties.pdf, PDF page 100).
  % SOURCE QUOTE:
  % "For any variety \(V\), the symmetric group \(S_r\) on \(r\) letters acts on
  %  the product of \(r\) copies \(V^r\) of \(V\) by permuting the factors, and
  %  we want to define the \(r\)th symmetric power \(V^{(r)}\) of \(V\) to be the
  %  quotient \(S_r \backslash V^r\).  The next proposition demonstrates the
  %  existence of \(V^{(r)}\) and lists its main properties.
  %    A morphism \(\varphi\colon V^r\to T\) is said to be {\bf symmetric} if
  %  \(\varphi\sigma=\varphi\) for all \(\sigma\) in \(S_r\).
  %  PROPOSITION 3.1. Let \(V\) be a variety over \(k\).  Then there exists a
  %  variety \(V^{(r)}\) and a symmetric morphism \(\pi\colon V^r\to V^{(r)}\)
  %  having the following properties:
  %    (a) as a topological space, \((V^{(r)},\pi)\) is the quotient of \(V^r\)
  %        by \(S_r\);
  %    (b) for any open affine subset \(U\) of \(V\), \(U^{(r)}\) is an open
  %        affine subset of \(V^{(r)}\) and
  %             $\Gamma(U^{(r)},\mathcal O_{V^{(r)}}) = \Gamma(U^r,
  %             \mathcal O_{V^r})^{S_r}$
  %        (set of elements fixed by the action of \(S_r\)).
  %  The pair \((V^{(r)},\pi)\) has the following universal property: every
  %  symmetric \(k\)-morphism \(\varphi\colon V^r\to T\) factors uniquely
  %  through \(\pi\).
  %  The map \(\pi\) is finite, surjective, and separable."
  \textit{Source: Milne, \emph{Abelian Varieties}, Proposition~3.1, \S III.3, p.~94.}

  Let \(C\) be a complete nonsingular curve over an algebraically closed field
  \(\bar k\) and let \(g \ge 1\). The \emph{\(g\)-th symmetric power}
  \(\mathrm{Sym}^g C\) is a \(\bar k\)-scheme equipped with a finite surjective
  symmetric morphism \(\pi \colon C^g \to \mathrm{Sym}^g C\) characterised by
  the following universal property: for every \(\bar k\)-scheme \(T\) and every
  symmetric \(\bar k\)-morphism \(\varphi \colon C^g \to T\) (i.e.\ a morphism
  satisfying \(\varphi \circ \sigma = \varphi\) for all \(\sigma \in S_g\)
  acting by permutation of the factors of \(C^g\)) there exists a unique
  morphism \(\overline\varphi \colon \mathrm{Sym}^g C \to T\) with
  \(\overline\varphi \circ \pi = \varphi\). As a topological space
  \(\mathrm{Sym}^g C\) is the quotient of \(C^g\) by \(S_g\), and for every open
  affine \(U \subseteq C\) the open subset \(U^{(g)} \subseteq \mathrm{Sym}^g
  C\) is affine with global sections the \(S_g\)-invariants
  \(\Gamma(U^g, \mathcal O_{C^g})^{S_g}\).
\end{definition}

\medskip

\textbf{Mathlib status.} Mathlib has no formalised \(g\)-th symmetric power of
a scheme: only \texttt{Sym}/\texttt{SymmetricPower} for types and modules.
The corresponding scheme-theoretic construction must be supplied
project-side. The intended Lean construction follows Milne's recipe of
\cref{def:symmetric_power_curve}: on affine open \(U = \Spec A \subseteq C\)
form \(\Spec(A^{\otimes g})^{S_g}\), then glue via the standard affine
patching for an open cover \(C = \bigcup U_i\). This sub-build is a
non-trivial dependency of A.4.d that the prover lane must absorb before the
main universal property closes; a \texttt{\textbackslash notready} marker is
appropriate on the corresponding Lean target until the construction lands.
For the proof of Proposition~3.1 itself, Milne refers to Mumford 1970,
\S II~\S 7, p.~66 and \S III \S 11, p.~112, for the patching details (the
quoted reference at the end of Milne's Proof on \S III.3, p.~94).

\section{The morphism \(\mathrm{Sym}^g \varphi \colon \mathrm{Sym}^g C \to A\)}
\label{sec:albaneseup_sym_g_av_map}

The second sub-lemma is the elementary fact that the symmetric assignment
\((P_1, \ldots, P_g) \mapsto \varphi(P_1) + \cdots + \varphi(P_g)\) from
\(C^g\) to \(A\) descends through \(\pi \colon C^g \to \mathrm{Sym}^g C\). The
key point is that the addition morphism \(A^g \to A\) is itself
\(S_g\)-symmetric because \(A\) is commutative as a group scheme.

\begin{lemma}\leanok
[Symmetrisation of a curve-to-AV morphism]
  \label{lem:symmetric_product_av_map}
  \lean{AlgebraicGeometry.Pic0.symmetricPowerAVMap}
  \uses{def:symmetric_power_curve}

  % SOURCE: [Milne, Abelian Varieties], Proof of Proposition~6.1, \S III.6,
  % doc p.~104 (read from references/abelian-varieties.pdf, PDF page 110).
  % SOURCE QUOTE:
  % "PROOF.  Consider the map
  %     \((P_1,...,P_g)\mapsto\sum_i\psi(P_i)\colon C^g\to A.\)
  %  Clearly this is symmetric, and so it factors through \(C^{(g)}\)."
  \textit{Source: Milne, \emph{Abelian Varieties}, Proof of Proposition~6.1,
  \S III.6, p.~104.}

  Let \(\varphi \colon C \to A\) be a morphism into an abelian variety \(A\)
  over \(\bar k\). The composition
  \[
    \mathrm{add}_A \circ \varphi^{\times g}
    \;\colon\; C^g \;\longrightarrow\; A^g \;\longrightarrow\; A,
    \qquad (P_1, \ldots, P_g) \;\longmapsto\; \varphi(P_1) + \cdots + \varphi(P_g),
  \]
  is \(S_g\)-equivariant (where \(S_g\) acts on \(C^g\) by permutation of the
  factors and on \(A\) trivially). Consequently, by the universal property of
  \(\mathrm{Sym}^g C\) (\cref{def:symmetric_power_curve}), there is a unique
  morphism of \(\bar k\)-schemes
  \[
    \mathrm{Sym}^g \varphi \;\colon\; \mathrm{Sym}^g C \;\longrightarrow\; A
  \]
  such that
  \(\mathrm{Sym}^g \varphi \circ \pi = \mathrm{add}_A \circ \varphi^{\times g}\).
\end{lemma}

\begin{proof}
  
  Since \(A\) is a commutative group scheme, the \(g\)-fold addition morphism
  \(\mathrm{add}_A \colon A^g \to A\) satisfies \(\mathrm{add}_A \circ \tau =
  \mathrm{add}_A\) for every \(\tau \in S_g\) permuting the factors of \(A^g\)
  (this is the scheme-theoretic statement that summation in a commutative
  monoid is independent of the order of summands; on \(T\)-points it is the
  identity \(a_1 + \cdots + a_g = a_{\tau(1)} + \cdots + a_{\tau(g)}\) in the
  abelian group \(A(T)\)). The diagonal action of \(\tau\) on \(\varphi^{\times
  g} \colon C^g \to A^g\) is also given by permutation of factors. Therefore
  \[
    (\mathrm{add}_A \circ \varphi^{\times g}) \circ \tau
    \;=\; \mathrm{add}_A \circ (\varphi^{\times g} \circ \tau)
    \;=\; \mathrm{add}_A \circ (\tau \circ \varphi^{\times g})
    \;=\; (\mathrm{add}_A \circ \tau) \circ \varphi^{\times g}
    \;=\; \mathrm{add}_A \circ \varphi^{\times g},
  \]
  so the composition \(\mathrm{add}_A \circ \varphi^{\times g}\) is symmetric
  in the sense of \cref{def:symmetric_power_curve}. The universal property
  of \(\mathrm{Sym}^g C\) produces the unique factorisation
  \(\mathrm{Sym}^g \varphi \colon \mathrm{Sym}^g C \to A\).
\end{proof}

\section{The birational morphism \(\mathrm{Sym}^g C \dashrightarrow J\)}
\label{sec:albaneseup_sym_g_to_jacobian}

The third sub-lemma is the birational identification of \(\mathrm{Sym}^g C\)
with \(J\). After choosing the basepoint \(P_0 \in C(\bar k)\), the cycle map
\((P_1, \ldots, P_g) \mapsto [\mathcal O_C(P_1 + \cdots + P_g - g P_0)]\)
defines a morphism \(f^{(g)} \colon \mathrm{Sym}^g C \to J\) which is
birational by Milne III~Theorem~5.1(a).

\begin{lemma}\leanok
[Birational morphism from \(\mathrm{Sym}^g C\) to \(J\)]
  \label{lem:symmetric_product_to_jacobian}
  \lean{AlgebraicGeometry.Pic0.symmetricPowerToJacobian}
  \uses{lem:abel_jacobi_morphism, def:symmetric_power_curve, def:jacobian_scheme, lem:jacobian_scheme_grpObj, lem:jacobian_scheme_isProper, lem:jacobian_scheme_smooth, lem:jacobian_scheme_geomIrred}

  % SOURCE: [Milne, Abelian Varieties], Theorem~5.1(a), \S III.5,
  % doc p.~101 (read from references/abelian-varieties.pdf, PDF page 107).
  % SOURCE QUOTE:
  % "Let \(f^r\) be the map \(C^r\to J\) sending \((P_1,...,P_r)\) to
  %  \(f(P_1)+\cdots+f(P_r)\).  On points, \(f^r\) is the map
  %  \((P_1,...,P_r)\mapsto[P_1+\cdots+P_r-rP]\).  Clearly it is symmetric,
  %  and so induces a map \(f^{(r)}\colon C^{(r)}\to J\).  We can regard
  %  \(f^{(r)}\) as being the map sending an effective divisor \(D\) of degree
  %  \(r\) on \(C\) to the linear equivalence class of \(D-rP\).  The fibre of
  %  the map \(f^{(r)}\colon C^{(r)}(k)\to J(k)\) containing \(D\) can be
  %  identified with the space of effective divisors linearly equivalent
  %  to \(D\), that is, with the linear system \(|D|\).  The image of \(C^{(r)}\)
  %  in \(J\) is a closed subvariety \(W^r\) of \(J\), which can also be written
  %  \(W^r=f(C)+\cdots+f(C)\) (\(r\) summands).
  %  THEOREM 5.1.  (a) For all \(r\le g\), the morphism
  %  \(f^{(r)}\colon C^{(r)}\to W^r\) is birational; in particular,
  %  \(f^{(g)}\) is a birational map from \(C^{(g)}\) onto \(J\)."
  \textit{Source: Milne, \emph{Abelian Varieties}, Theorem~5.1(a) and the
  preceding paragraph, \S III.5, p.~101.}

  Let \(C\) be a complete nonsingular curve of genus \(g > 0\) over \(\bar k\),
  let \(P_0 \in C(\bar k)\) be a basepoint, and write \(J := \Pic^0_{C/k}\),
  \(\iota_{P_0} \colon C \to J\) for the Abel--Jacobi morphism of
  \cref{lem:abel_jacobi_morphism}. The regular morphism
  \(\iota_{P_0}^{\times g} \colon C^g \to J^g\) composed with the \(g\)-fold
  group-law addition \(\mathrm{add}_J \colon J^g \to J\) descends through
  \(\pi \colon C^g \to \mathrm{Sym}^g C\) to a regular morphism
  \[
    f^{(g)} \;\colon\; \mathrm{Sym}^g C \;\longrightarrow\; J,
    \qquad (P_1, \ldots, P_g) \;\longmapsto\;
      \bigl[\mathcal O_C(P_1 + \cdots + P_g - g P_0)\bigr].
  \]
  The morphism \(f^{(g)}\) is birational: it is dominant (its image is
  \(J\), the \(g\)-fold sum \(\iota_{P_0}(C) + \cdots + \iota_{P_0}(C)\) exhausts
  \(J\) since \(\dim J = g\) and the iterated sum reaches full dimension), and
  its generic fibre is a single \(\bar k\)-point (a degree-\(g\) effective
  divisor \(D\) on \(C\) with \(h^0(D) = 1\) is the unique element of its linear
  system \(|D|\); such divisors form a dense open subset of \(\mathrm{Sym}^g
  C\) -- a consequence of \(\dim |D| = 0\) generically on a curve of genus \(g\)
  by Riemann--Roch). Equivalently, \(f^{(g)}\) is a birational morphism in
  the sense that there exist dense opens
  \(U \subseteq \mathrm{Sym}^g C\), \(V \subseteq J\) such that
  \(f^{(g)}|_U \colon U \to V\) is an isomorphism.
\end{lemma}

% NOTE (iter-199): the birationality proof below derives the generic-fibre
% identification from Riemann--Roch at \(\deg D = g\) (specifically
% \(h^0(D) - h^1(D) = 1\) generically). Riemann--Roch is the Route C
% chain (\texttt{RiemannRoch_RRFormula.tex} et al.), currently PAUSED
% per USER 2026-05-28 standing directive. This lemma is therefore
% gated on Route C re-engagement (or on a parallel construction that
% sidesteps RR). Flagged here so a future prover does not attempt to
% formalize the birationality argument until the RR substrate is
% available.

\begin{proof}

  Symmetry of \(\mathrm{add}_J \circ \iota_{P_0}^{\times g}\) is the
  specialisation of \cref{lem:symmetric_product_av_map} to the abelian
  variety \(A := J\) and the morphism \(\varphi := \iota_{P_0}\); the universal
  property of \(\mathrm{Sym}^g C\) produces the regular morphism \(f^{(g)}\).
  The closed-form description on \(\bar k\)-points,
  \(f^{(g)}(P_1, \ldots, P_g) = [\mathcal O_C(P_1 + \cdots + P_g - g P_0)]\),
  follows because \(\iota_{P_0}(P_i) = [\mathcal O_C(P_i - P_0)]\) by
  \cref{lem:abel_jacobi_morphism} and the group law on \(J\) is the
  tensor-product of degree-zero line bundles.

  We verify birationality. The image of \(f^{(g)}\) in \(J\) is the iterated
  sum \(\iota_{P_0}(C) + \cdots + \iota_{P_0}(C)\) (\(g\) summands), which is
  closed (it is the image of the proper morphism
  \(\mathrm{Sym}^g C \to J\), \(\mathrm{Sym}^g C\) being proper because \(C^g\)
  is proper and \(\pi\) is finite surjective). By the dimension count
  \(\dim \mathrm{Sym}^g C = g = \dim J\) (Riemann--Roch and the standing
  hypothesis \(g > 0\)), the image \(W^g \subseteq J\) has dimension at most
  \(g\), and equality forces \(W^g = J\) since \(J\) is irreducible. Hence
  \(f^{(g)}\) is surjective with \(\dim\)-equal source and target.

  For the generic-fibre identification, Riemann--Roch
  \(h^0(D) - h^1(D) = \deg(D) + 1 - g\) at \(\deg D = g\) gives \(h^0(D) - h^1(D)
  = 1\), so for a general \(D \in \mathrm{Sym}^g C(\bar k)\) with \(h^1(D) =
  0\) (the locus where \(h^1\) jumps is closed of positive codimension on a
  smooth curve, by upper semicontinuity of cohomology), we have \(h^0(D) =
  1\) and the linear system \(|D|\) consists of \(D\) alone. Therefore the
  fibre of \(f^{(g)}\) over the generic point of \(J\) contains a single point
  \(D\), i.e.\ \(f^{(g)}\) is birational onto \(J\).

  By \cref{def:pic_scheme} the underlying \(\bar k\)-scheme \(J\) is the
  Picard scheme \(\Pic^0_{C/k}\) obtained at A.3, so the moduli-theoretic
  reading of \(f^{(g)}\) (effective divisor class \(\mapsto\) degree-zero
  line bundle class via \(D \mapsto [\mathcal O_C(D - g P_0)]\)) is the
  same as the Yoneda image of the \(\mathrm{Sym}^g C\)-point
  \([\mathcal O_{C \times \mathrm{Sym}^g C}(\mathcal D_g - g\{P_0\} \times
  \mathrm{Sym}^g C)]\), where \(\mathcal D_g \subseteq C \times \mathrm{Sym}^g
  C\) is the universal effective Cartier divisor of degree \(g\) on \(C\). We
  do not need this Yoneda repackaging for the proof of the Albanese UP --
  the symmetric-power route uses only the morphism \(f^{(g)}\) and its
  birationality.
\end{proof}

\section{Descent through the birational quotient \(\sigma\)}
\label{sec:albaneseup_descent}

The final auxiliary lemma packages the two-step descent that promotes a
morphism on \(C^g\) (after passing through \(\mathrm{Sym}^g C\) via
\cref{lem:symmetric_product_av_map} and inverting \(f^{(g)}\) via
\cref{lem:symmetric_product_to_jacobian}) to a regular morphism on
\(J\). The output is a rational map \(J \dashrightarrow A\), which is then
upgraded to a regular morphism via Milne's Theorem~3.2
(\cref{thm:rational_map_to_av_extends}). This is the unique step of the
proof that consumes A.4.c.

\begin{lemma}\leanok
[Descent of \(\mathrm{Sym}^g \varphi\) to a regular morphism \(J \to A\)]
  \label{lem:descent_through_birational_sigma}
  \lean{AlgebraicGeometry.Pic0.descentThroughBirationalSigma}
  \uses{lem:symmetric_product_av_map, lem:symmetric_product_to_jacobian,
        thm:rational_map_to_av_extends}

  % SOURCE: [Milne, Abelian Varieties], Proof of Proposition~6.1, \S III.6,
  % doc p.~104 (read from references/abelian-varieties.pdf, PDF page 110).
  % SOURCE QUOTE:
  % "It therefore defines a rational map \(\psi\colon J\to A\), which (I 3.2)
  %  shows to be a regular map."
  \textit{Source: Milne, \emph{Abelian Varieties}, Proof of Proposition~6.1,
  \S III.6, p.~104 (the sentence ``which (I~3.2) shows to be a regular
  map'').}

  Let \(\varphi \colon C \to A\) be a morphism into an abelian variety \(A\)
  over \(\bar k\). Let \(\mathrm{Sym}^g \varphi \colon \mathrm{Sym}^g C \to A\)
  be the symmetrised morphism of \cref{lem:symmetric_product_av_map} and
  \(f^{(g)} \colon \mathrm{Sym}^g C \to J\) the birational morphism of
  \cref{lem:symmetric_product_to_jacobian}. Then there is a unique
  regular morphism
  \[
    \psi \;\colon\; J \;\longrightarrow\; A
  \]
  such that \(\psi \circ f^{(g)} = \mathrm{Sym}^g \varphi\) on the dense open
  on which \(f^{(g)}\) is an isomorphism (equivalently, \(\psi\) is the unique
  regular morphism extending the rational map
  \(\mathrm{Sym}^g \varphi \circ (f^{(g)})^{-1} \colon J \dashrightarrow
  A\)).
\end{lemma}

\begin{proof}
  
  By \cref{lem:symmetric_product_to_jacobian}, there exist dense opens
  \(U \subseteq \mathrm{Sym}^g C\) and \(V \subseteq J\) with \(f^{(g)}|_U \colon
  U \xrightarrow{\sim} V\) an isomorphism. Setting \(\psi_0 \colon V \to A\)
  to be the composition \(\mathrm{Sym}^g \varphi|_U \circ (f^{(g)}|_U)^{-1}
  \colon V \to U \to A\) produces a rational map \(J \dashrightarrow A\).

  The target \(A\) is an abelian variety and the source \(J\) is smooth over
  \(\bar k\) (it is itself an abelian variety, hence nonsingular). Milne's
  Theorem~3.2 (\cref{thm:rational_map_to_av_extends}) therefore promotes
  \(\psi_0\) to a unique regular morphism \(\psi \colon J \to A\) extending
  \(\psi_0\) on \(V\). Uniqueness of \(\psi\) follows from the
  reduced-and-separated agreement principle: two regular morphisms
  \(J \to A\) agreeing on a dense open are equal.
\end{proof}

\section{Connecting the two factorisation forms}
\label{sec:albaneseup_connecting}

The headline universal property is phrased in the \emph{Albanese form}
\(\varphi = \psi \circ \iota_{P_0}\), whereas the descent lemma
\cref{lem:descent_through_birational_sigma} delivers the
\emph{symmetric-power form}
\(\mathrm{Sym}^g \varphi = \psi \circ f^{(g)}\). The following biconditional
identifies the two forms for every candidate \(\psi\), so that the existence
and uniqueness statement of \cref{thm:albanese_universal_property} reduces to
that of the descent lemma.

\begin{lemma}\leanok
[Albanese form \(\Leftrightarrow\) symmetric-power form]
  \label{lem:albanese_eq_iff_symmetricPower_eq}
  \lean{AlgebraicGeometry.Pic0.albanese_eq_iff_symmetricPower_eq}
  \uses{lem:abel_jacobi_morphism, lem:symmetric_product_av_map, lem:symmetric_product_to_jacobian}

  Let \(\varphi \colon C \to A\) be a pointed morphism into an abelian variety
  \(A\) over \(\bar k\) with \(\varphi(P_0) = \eta_A\). For every candidate
  morphism \(\psi \colon J \to A\), the Albanese-form factorisation
  \(\varphi = \psi \circ \iota_{P_0}\) holds if and only if the
  symmetric-power-form factorisation
  \(\mathrm{Sym}^g \varphi = \psi \circ f^{(g)}\) holds, where
  \(\iota_{P_0}\) is the Abel--Jacobi morphism of
  \cref{lem:abel_jacobi_morphism}, \(\mathrm{Sym}^g \varphi\) the symmetrised
  morphism of \cref{lem:symmetric_product_av_map}, and \(f^{(g)}\) the
  birational morphism of \cref{lem:symmetric_product_to_jacobian}. The forward
  direction precomposes the Albanese equation with the symmetric \(g\)-fold sum
  and descends through \(\pi \colon C^g \to \mathrm{Sym}^g C\); the reverse
  direction restricts the symmetric-power equation along the embedding
  \(Q \mapsto (Q, P_0, \ldots, P_0)\) and uses \(\varphi(P_0) = \eta_A\).
\end{lemma}

\begin{proof} Proved directly in Lean. \end{proof}

\section{Proof of the universal property}
\label{sec:albaneseup_proof}

% SOURCE QUOTE PROOF:
% "PROOF.  Consider the map
%    \(\)(P_1,\ldots,P_g)\mapsto\sum_i\psi(P_i)\colon C^g\to A.\(\)
%  Clearly this is symmetric, and so it factors through \(C^{(g)}\).  It
%  therefore defines a rational map \(\psi\colon J\to A\), which (I 3.2)
%  shows to be a regular map.  It is clear from the construction that
%  \(\psi\circ f^P=\varphi\) (note that \(f^P\) is the composite of $Q\mapsto
%  Q+(g-1)P\colon C\to C^{(g)}\( with \)f^{(g)}\colon C^{(g)}\to J$).  In
%  particular, \(\psi\) maps 0 to 0, and (I 1.2) shows that it is therefore
%  a homomorphism.  If \(\psi'\) is a second homomorphism such that
%  \(\psi'\circ f^P=\varphi\), then \(\psi\) and \(\psi'\) agree on
%  \(f^P(C)+\cdots+f^P(C)\) (\(g\) copies), which is the whole of \(J\)."
\begin{proof}[Proof of \cref{thm:albanese_universal_property}]
  
  The proof follows Milne's Proposition~6.1 verbatim, using the four
  sub-lemmas of this chapter:
  \cref{lem:abel_jacobi_morphism} (Abel--Jacobi morphism \(\iota_{P_0}\)),
  \cref{lem:symmetric_product_av_map} (symmetrisation of \(\varphi\) to
  \(\mathrm{Sym}^g \varphi \colon \mathrm{Sym}^g C \to A\)),
  \cref{lem:symmetric_product_to_jacobian} (birational morphism
  \(f^{(g)} \colon \mathrm{Sym}^g C \to J\)), and
  \cref{lem:descent_through_birational_sigma} (descent of \(\mathrm{Sym}^g
  \varphi\) to a regular morphism \(\psi \colon J \to A\) via
  \cref{thm:rational_map_to_av_extends}).
  \begin{enumerate}
    \item \emph{Symmetrisation.} Consider the map \(C^g \to A\) sending
      \((P_1, \ldots, P_g) \mapsto \varphi(P_1) + \cdots + \varphi(P_g)\)
      (Milne's typeset proof writes \(\sum_i \psi(P_i)\) in this line, but
      the construction is the symmetrisation of the input \(\varphi\); we
      use the unambiguous symbol \(\varphi\) throughout). Since \(A\) is
      commutative this map is \(S_g\)-symmetric. By
      \cref{lem:symmetric_product_av_map} it factors through the
      symmetric power \(\mathrm{Sym}^g C\), defining a regular morphism
      \[
        \widetilde \psi \;:=\; \mathrm{Sym}^g \varphi
        \;\colon\; \mathrm{Sym}^g C \;\longrightarrow\; A.
      \]
    \item \emph{Descent to a regular morphism on \(J\).} By
      \cref{lem:symmetric_product_to_jacobian}, the morphism
      \(f^{(g)} \colon \mathrm{Sym}^g C \to J\) is birational. The
      composite \(\widetilde \psi \circ (f^{(g)})^{-1}\) is, a priori,
      only a rational map \(J \dashrightarrow A\). By
      \cref{lem:descent_through_birational_sigma} -- whose proof is an
      application of \cref{thm:rational_map_to_av_extends} (Milne
      Theorem~I~3.2, rational map from a nonsingular variety to an
      abelian variety is regular) -- this rational map extends uniquely
      to a regular morphism
      \[
        \psi \;\colon\; J \;\longrightarrow\; A.
      \]
    \item \emph{Factorisation \(\psi \circ \iota_{P_0} = \varphi\).}
      Milne's identification: \(\iota_{P_0}\) is the composite
      \[
        C \;\xrightarrow{\;Q \mapsto Q + (g-1) P_0\;}\;
        \mathrm{Sym}^g C \;\xrightarrow{\;f^{(g)}\;}\; J,
      \]
      where the first arrow sends a point \(Q \in C(\bar k)\) to the
      effective divisor \(Q + (g-1) P_0 \in \mathrm{Sym}^g C(\bar k)\)
      (its image under \(f^{(g)}\) in \(J\) is
      \([\mathcal O_C(Q + (g-1) P_0 - g P_0)] = [\mathcal O_C(Q - P_0)]
      = \iota_{P_0}(Q)\), matching \cref{lem:abel_jacobi_morphism}). On
      \(\bar k\)-points,
      \[
        \psi(\iota_{P_0}(Q)) = \widetilde\psi(Q + (g-1) P_0)
          = \varphi(Q) + (g-1)\varphi(P_0) = \varphi(Q),
      \]
      using \(\varphi(P_0) = \eta_A\) and that \(\eta_A\) is the identity of
      \(A\). Hence \(\psi \circ \iota_{P_0} = \varphi\) on \(\bar k\)-points,
      and therefore \(\psi \circ \iota_{P_0} = \varphi\) as morphisms of
      \(\bar k\)-schemes (both sides are regular morphisms \(C \to A\) from
      a reduced separated scheme to a separated scheme, agreeing on
      dense \(\bar k\)-points, hence agreeing everywhere by the
      reduced-and-separated agreement principle).
    \item \emph{Homomorphism property.} The morphism \(\psi\) sends
      \(\eta_J\) to \(\eta_A\): indeed,
      \(\psi(\eta_J) = \psi(\iota_{P_0}(P_0)) = \varphi(P_0) = \eta_A\).
      A regular morphism of abelian varieties sending the identity to
      the identity is a homomorphism of group schemes, by Milne
      Corollary~I~1.2 (a consequence of the Rigidity Lemma
      \cref{thm:rigidity_lemma}, see
      \cref{chap:AbelianVarietyRigidity}, Corollary~1.2). Therefore
      \(\psi\) is a homomorphism of group schemes \(J \to A\).
    \item \emph{Uniqueness.} Suppose \(\psi'\) is a second homomorphism
      with \(\psi' \circ \iota_{P_0} = \varphi\). Then \(\psi\) and
      \(\psi'\) agree on the set-theoretic image \(\iota_{P_0}(C) \subseteq
      J\). Both are homomorphisms, so they agree on the subgroup
      generated by \(\iota_{P_0}(C)\) -- which contains the \(g\)-fold sum
      \(\iota_{P_0}(C) + \cdots + \iota_{P_0}(C)\) (\(g\) copies). The
      \(g\)-fold sum is the image of
      \(\mathrm{add}_J \circ \iota_{P_0}^{\times g} \colon C^g \to J\),
      which descends to \(f^{(g)} \colon \mathrm{Sym}^g C \to J\), which
      is surjective by \cref{lem:symmetric_product_to_jacobian}.
      Therefore \(\psi\) and \(\psi'\) agree on a dense open subset of \(J\)
      (in fact on all of \(J\) as a set), hence everywhere by the
      reduced-and-separated agreement principle.
  \end{enumerate}
  This completes the proof.
\end{proof}

\section{Wiring for the positive-genus Albanese witness}
\label{sec:albaneseup_witness_wiring}

\cref{thm:albanese_universal_property} supplies the universal-factorisation
field \(\texttt{isAlbaneseFor}\) of an Albanese witness for \((C, P_0)\) in the
sense of \cref{def:JacobianWitness} (\cref{chap:Jacobian}); the remaining
fields are supplied by A.3:
\begin{itemize}
  \item \(J := \Pic^0_{C/k}\) (the underlying scheme).
  \item \(\texttt{grpObj}\): the group-object structure on \(J\) obtained from the
    abelian-group structure on \(\Pic^0_{C/k}(T)\) for varying \(T\), refined to a
    group object in the category of \(\bar k\)-schemes (A.3, A.2.c group-scheme
    refinement).
  \item \(\texttt{proper}\): properness of \(\Pic^0_{C/k} \to \Spec \bar k\) from
    the closed-immersion property of the identity component in the proper
    representing Picard scheme (A.3; cf.\
    \cref{thm:Jacobian_proper}'s discussion).
  \item \(\texttt{smoothGenus}\): smoothness of relative dimension \(g(C)\) from
    the deformation-theoretic tangent-space identification with \(H^1(C,
    \mathcal O_C)\), which has dimension \(g(C)\) by Riemann--Roch.
  \item \(\texttt{geomIrred}\): geometric irreducibility of the identity
    component of a \(\bar k\)-group scheme locally of finite type (Kleiman,
    \emph{The Picard scheme}, the identity-component lemma and its
    \(\Pic^0_{C/k}\) specialisation; cf.\
    \cref{thm:Jacobian_geomIrred}'s discussion).
  \item \(\texttt{isAlbaneseFor}\): the present
    \cref{thm:albanese_universal_property}, applied at each
    \(P_0 \in C(\bar k)\).
\end{itemize}

The witness so produced is over \(\bar k\). The Galois-descent step from \(\bar
k\) back to the original base field \(k\) -- packaged at
\cref{chap:Jacobian}, sub-step C.2.f -- closes the loop and supplies the
\(k\)-side Albanese witness consumed by \cref{def:Jacobian}.

\section{Lean encoding}
\label{sec:albaneseup_lean_encoding}

The Lean target file is the new
\texttt{AlgebraicJacobian/Albanese/AlbaneseUP.lean} (declared at the top of
this chapter). The five target declarations are:
\begin{itemize}
  \item \texttt{AlgebraicGeometry.Pic0.abelJacobi}: the regular morphism
    \(\iota_{P_0} \colon C \to \Pic^0_{C/k}\) of \cref{lem:abel_jacobi_morphism}.
    The signature exposes the marked point \(P_0\) as an explicit argument and
    the Picard-scheme construction \(\Pic^0_{C/k}\) as an existing object.
  \item \texttt{AlgebraicGeometry.Pic0.SymmetricPower}: the scheme-theoretic
    \(g\)-th symmetric power \(\mathrm{Sym}^g C\) of \cref{def:symmetric_power_curve},
    packaged with the canonical projection
    \(\pi \colon C^g \to \mathrm{Sym}^g C\) and its universal property
    (factorisation of every \(S_g\)-symmetric morphism out of \(C^g\)). \emph{This
    declaration has no Mathlib analogue}: Mathlib's \texttt{Sym} /
    \texttt{SymmetricPower} cover types and modules only; the
    scheme-theoretic construction must be supplied project-side, following
    Milne's affine-and-glue recipe of \cref{def:symmetric_power_curve}.
  \item \texttt{AlgebraicGeometry.Pic0.symmetricPowerAVMap}: the
    symmetrised morphism
    \(\mathrm{Sym}^g \varphi \colon \mathrm{Sym}^g C \to A\) produced by
    \cref{lem:symmetric_product_av_map} from an input morphism
    \(\varphi \colon C \to A\) into an abelian variety.
  \item \texttt{AlgebraicGeometry.Pic0.symmetricPowerToJacobian}: the
    birational morphism \(f^{(g)} \colon \mathrm{Sym}^g C \to J\) of
    \cref{lem:symmetric_product_to_jacobian}, together with the data of a
    dense open isomorphism \(U \xrightarrow{\sim} V\) realising
    birationality. The signature exposes the basepoint \(P_0\) and consumes
    \texttt{AlgebraicGeometry.Pic0.abelJacobi}.
  \item \texttt{AlgebraicGeometry.Pic0.descentThroughBirationalSigma}: the
    descent lemma of \cref{lem:descent_through_birational_sigma},
    consuming \texttt{symmetricPowerAVMap} and
    \texttt{symmetricPowerToJacobian} together with
    \texttt{Albanese.Thm32RationalMapExtension.extend\_to\_av}
    (\cref{thm:rational_map_to_av_extends}) to produce the regular
    morphism \(\psi \colon J \to A\).
  \item \texttt{AlgebraicGeometry.Pic0.albanese\_universal\_property}: the
    main theorem, the existence-and-uniqueness statement of
    \cref{thm:albanese_universal_property}. The signature is
    \[
      \texttt{P\_0 : C(\bar k)} \;\to\;
      \texttt{A : AbelianVariety \(\bar k\)} \;\to\;
      (\varphi : C \to A) \;\to\;
      (h\varphi : \varphi \circ P_0 = \eta_A) \;\to\;
      \exists ! \, \psi \colon J \to A, \;\; \psi \circ \iota_{P_0} = \varphi
      \;\wedge\; \psi \text{ is a group homomorphism}.
    \]
\end{itemize}

The proof body of \texttt{albanese\_universal\_property} consumes its four
sub-lemma exports plus \texttt{Albanese.Thm32RationalMapExtension.extend\_to\_av}
(\cref{thm:rational_map_to_av_extends}) and Milne's Corollary~I~1.2
(\cref{thm:rigidity_lemma} consequence, registered at
\cref{chap:AbelianVarietyRigidity}); the wiring follows the five-step proof
of \cref{thm:albanese_universal_property}.

The Picard-scheme inputs (\cref{def:pic_scheme} of
\cref{chap:Picard_FGAPicRepresentability}) are imported from
\texttt{AlgebraicJacobian/Picard/FGAPicRepresentability.lean} (when that file
materialises along the A.2.c row) and the identity-component / degree-map
refinement is gated on the A.3 row (which does not yet have a dedicated
chapter; see ``Notes for Plan Agent'' below). At the time of writing, the
prover-side declaration body
\texttt{AlgebraicGeometry.Pic0.albane\-se\_universal\_property} should be
skeletoned as a \texttt{sorry} that consumes the A.3 instances as
hypotheses with named placeholders, and \texttt{Sym\-met\-ric\-Power} should
be skeletoned with a \texttt{notready} marker until the affine-and-glue
sub-build lands. The remaining sub-lemma bodies become routine once both
inputs are available.

\section{Out of scope}
\label{sec:albaneseup_out_of_scope}

This chapter packages \emph{only} the Albanese universal property
(\cref{thm:albanese_universal_property} = Milne Proposition~III~6.1) and the
four internal sub-lemmas of the symmetric-power proof
(\cref{lem:abel_jacobi_morphism},
\cref{def:symmetric_power_curve},
\cref{lem:symmetric_product_av_map},
\cref{lem:symmetric_product_to_jacobian},
\cref{lem:descent_through_birational_sigma}), plus the wiring of these into
a positive-genus Albanese witness. The following adjacent content is
explicitly out of scope:
\begin{itemize}
  \item \textbf{Representability of \(\Pic_{C/k}\) (A.2.c).} Proved in
    \cref{chap:Picard_FGAPicRepresentability} as
    \cref{thm:fga_pic_representability}; this chapter only consumes it.
  \item \textbf{The identity component \(\Pic^0_{C/k}\) and the degree map
    (A.3).} The identity component refinement of the Picard scheme together
    with its degree map is the subject of a forthcoming chapter; this chapter
    consumes \(\Pic^0_{C/k}\) as a smooth proper geometrically irreducible
    \(\bar k\)-group scheme of dimension \(g\) but does not construct it.
  \item \textbf{Milne's Theorem~3.2 (A.4.c).} Proved in
    \cref{chap:Albanese_Thm32RationalMapExtension} as
    \cref{thm:rational_map_to_av_extends}; this chapter only consumes it
    inside \cref{lem:descent_through_birational_sigma} to promote the
    rational map \(J \dashrightarrow A\) to a regular morphism.
  \item \textbf{The scheme-theoretic construction of \(\mathrm{Sym}^g C\).}
    Mathlib has no \(g\)-th symmetric power of a scheme; the
    \texttt{SymmetricPower} target of \cref{def:symmetric_power_curve} is a
    project-side sub-build (affine \(S_g\)-invariants
    \(\Spec(A^{\otimes g})^{S_g}\) on each open affine \(\Spec A \subseteq C\),
    glued by the standard open-cover patching of Milne~III.3
    Proposition~3.1 / Mumford~1970 II.7 \& III.11). The construction itself
    is a non-trivial dependency of A.4.d; this chapter \emph{specifies} the
    target but does not develop its construction.
  \item \textbf{The autoduality \(J \cong J^\vee\) (Milne~III \S 6.6) and the
    moduli / Poincar\'e-bundle Albanese route.} The iter-174 draft of this
    chapter proposed a route via the Poincar\'e bundle of \(A\) and the
    autoduality of \(J\), which gates on the theorem of the cube. The cube
    was excised from the project in iter-163; the autoduality / moduli
    route is therefore demoted and replaced by the symmetric-power route
    of this chapter. A reference recording of the moduli wording is given
    in the ``Alternative route history'' \texttt{\%~NOTE} block at the end
    of this file.
  \item \textbf{Milne Proposition~III~6.4 / Remark~6.5 / Theorem~6.6} (the
    divisorial-correspondence universal property for \(C \times C \to A\),
    the identification of \((J, F)\) with the Albanese variety in the sense
    of Lang, and the autoduality \(J \cong J^\vee\)). Not needed by the
    symmetric-power route.
  \item \textbf{Weil's reconstruction of \(J\) from the rational law of
    composition on \(C^{(g)}\) (Milne~III \S 7).} The project does not need
    this alternative construction.
  \item \textbf{The genus-\(0\) case.} The standing hypothesis of this chapter
    is \(g(C) > 0\) (Milne's standing hypothesis of \S III.6). The genus-\(0\)
    case is not special: when \(g(C) = 0\) one has \(H^1(C, \mathcal O_C) = 0\),
    so \(\Pic^0_{C/k} = \Spec k\) automatically, and the genus-\(0\) witness is
    simply the degenerate instance of the same universal property (every pointed
    morphism \(C \to A\) into an abelian variety is then constant).
  \item \textbf{Generalisations to a non-algebraically-closed base \(k\).}
    The chapter works over \(\bar k\) throughout. The descent of the witness
    back to \(k\) via Galois descent of morphism equality is the subject of
    \cref{chap:Jacobian}, sub-step C.2.f.
\end{itemize}

% =========================================================================
% NOTE: Alternative route history (iter-174 moduli / autoduality wording).
% =========================================================================
%
% The iter-174 draft of this chapter proposed an alternative ``moduli''
% route via the Poincar\'e bundle of \(A\) and the autoduality
% \(J \cong J^\vee\). That route was retired in iter-175 because its
% autoduality input depends on the theorem of the cube, which was excised
% from the project in iter-163. The wording is recorded here for
% reference only; it is NOT part of the active dependency graph and is
% intentionally left without \label/\lean/\uses tags so it does not feed
% the leanblueprint dependency graph.
%
% ---- (Retired) Lemma: Poincar\'e-bundle pullback along $\varphi \times
% \mathrm{id}$ ----
% Let \(A\) be an abelian variety over \(\bar k\), equipped with its identity
% section \(\eta_A\), and let \(\varphi \colon C \to A\) be a pointed morphism
% with \(\varphi(P_0) = \eta_A\). Let \(\mathcal P_A\) be a Poincar\'e sheaf
% on \(A \times A^\vee\) normalised so that
% \(\eta_A^*\mathcal P_A \cong \mathcal O_{A^\vee}\). Then the pullback
% line bundle
%   \(\mathcal L_\varphi := (\varphi \times \mathrm{id}_{A^\vee})^*\mathcal P_A\)
% on \(C \times A^\vee\) is trivialised along \(\{P_0\} \times A^\vee\) and
% defines a class in \(\Pic^\sharp_{C/k}(A^\vee)\) via the relative Picard
% presheaf.
%
% ---- (Retired) Lemma: Moduli classifier of a trivialised line bundle ----
% Let \(T\) be a pointed \(\bar k\)-scheme with base-point \(t\). A line bundle
% \(\mathcal L\) on \(C \times T\) trivialised along \(\{P_0\} \times T\), with
% degree-zero geometric fibres, classifies under the moduli interpretation
% of \(J = \Pic^0_{C/k}\) a unique morphism \(h_{\mathcal L} \colon T \to J\)
% with \(h_{\mathcal L}(t) = \eta_J\).
%
% ---- (Retired) Moduli-route construction of \(\psi\) ----
% Apply the two retired lemmas at \(T := A^\vee\) with the input
% \((\mathcal P_A, \varphi)\) to obtain \(h_\varphi \colon A^\vee \to J\);
% dualise via the autoduality \(J \cong J^\vee\) and the canonical
% principal polarisation \(\varphi_{\mathcal L(\Theta)} \colon J \to J^\vee\)
% (Milne Theorem~III~6.6) to obtain \(\psi \colon J \to A\).
% The cube-dependence enters at the autoduality step: Milne~III~6.6 is
% proved using the seesaw theorem, which is a direct corollary of the
% theorem of the cube.
%
% End of retired wording.
% =========================================================================
